% !PRINTED 501
\weberpage{501}
\begin{center}
\S. 154.\\[4pt]
\gesperrt{Divisors of Groups. Subsidiary Groups and conjugate Groups.}
\end{center}

We saw in \S.\,152 that a function of $m$ independent variables always determines a permutation group of $m$ figures, but that this does not always hold when definite quantities are put in place of the independent variables. By the \gesperrt{Galois} theory, however, this difference is evened out, or at least reduced to its kernel. Now let $F(x) = 0$ be any equation of degree $m$ in a given field $\Omega$ with the mutually distinct roots $\alpha, \alpha_1, \alpha_2 \ldots \alpha_{m-1}$.

The normal field
\[
\Omega(\varrho) = \Omega(\alpha, \alpha_1, \alpha_2 \ldots \alpha_{m-1})
\]
we shall denote by $N$.

By the \gesperrt{Galois} group $P$ of the field $N$, or of the equation $F(x) = 0$, we may understand (by \S.\,146, 149) either the group of substitutions of the field $N$, or also the permutation group of the indices of the $\alpha$ isomorphic with it.

Let $p$ be the degree of $P$, and let us denote the operations of $P$ (be they now substitutions of $N$ or permutations of the $\alpha$) by
\weq{1}{P = \pi,\ \pi_1,\ \pi_2, \ldots \pi_{p-1}}
If the system $Q$ of the permutations contained in $P$
\weq{2}{\varkappa,\ \varkappa_1,\ \varkappa_2, \ldots \varkappa_{q-1}.}
% Source check: the faint terminal mark in (2) is the printed period.
forms a group by itself, that is, if any two of the elements of $Q$, compounded with one another, again yield an element of $Q$, then the group $Q$ is called a \gesperrt{divisor} or \gesperrt{factor of} $P$\footnote{Also called \gesperrt{subgroup}.}).

We first derive an important general theorem on the divisors of a group.

If the degree $q$ of the divisor is smaller than the degree $p$ of $P$, we take an element $\pi_1$ of $P$ not contained in $Q$ and form the system
\weq{3}{Q\,\pi_1 = \varkappa\,\pi_1,\ \varkappa_1\,\pi_1 \ldots \varkappa_{q-1}\,\pi_1,}
%CONT>>
